\documentclass[12pt,a4paper]{article}
\usepackage[brazil]{babel}
\usepackage{geometry}
\usepackage{setspace}
\usepackage{hyperref}
\geometry{left=3cm,right=2cm,top=3cm,bottom=2cm}
\onehalfspacing
\title{Modelagem computacional de interacoes fisico-quimicas para materiais energeticos seguros}
\author{Autor}
\date{2026}
\begin{document}
\maketitle
\begin{abstract}
Este trabalho propoe um estudo computacional e nao experimental sobre modelos matematicos de interacoes fisico-quimicas, com foco em materiais energeticos e limites de seguranca.
\end{abstract}
\textbf{Palavras-chave:} quimica computacional; fisica quantica; modelagem matematica; seguranca.
\section{Introducao}O estudo investiga como modelos matematicos e calculos de estrutura eletronica podem comparar interacoes em materiais energeticos sem manipular fontes radioativas.

\section{Problema e hipotese}Problema: quais descritores computacionais predizem estabilidade e transferencia de carga em sistemas candidatos? Hipotese: energia eletronica, gap e densidade de carga distinguem tendencias, mas nao substituem validacao experimental institucional.

\section{Objetivos}Objetivo geral: construir um workflow reprodutivel. Objetivos especificos: revisar fontes; gerar estruturas; calcular propriedades; quantificar incerteza; documentar limites e riscos.

\section{Revisao de literatura}Revisar artigos primarios sobre betavoltaica e carbono-14, metodos HF/DFT, baterias eletroquimicas e seguranca radiologica. Registrar DOI, acesso legal, metodo, amostra, resultado e limitacao.

\section{Metodologia}Usar Python, PySCF e, quando autorizado, Psi4. Fixar versoes, geometria, base, funcional, convergencia, unidades e sementes. Comparar calculos com referencia publicada e registrar erro absoluto e relativo.

\section{Resultados e discussao}Os resultados serao preenchidos apenas com saidas reproduzidas, tabela de parametros, incertezas e criterio de aprovacao previamente definido.

\section{Conclusao e limitacoes}Concluir somente quanto aos modelos computacionais. Nao inferir viabilidade de fabricacao radioativa a partir de simulacao. Limitar qualquer aplicacao a laboratorio habilitado e regulamentado.

\bibliographystyle{plain}\bibliography{refs}
\end{document}
